\documentclass{article}
\usepackage{/globals/code}
\usepackage{/globals/anki}
\ankiprefix{assembler}
\ankitopic{Assembler}
\ankishow
\begin{document}
\section{Assembler}
Die Assemblersprache ist die niedrigste generelle Programmiersprache; dessen Anweisungen sind identisch zu denen der Maschinensprache des Prozessors.

Befehle sind kurze Abkürzungen, die für ganze Worte stehen. Die Operanten sind entweder Zahlen, oder wenn sie in eckigen Klammern stehen, die Adresse einer Zahl. Der Text hinter einem Semicolon ist ein Kommentar.
\begin{slsourcecode}{asm}
00 LDA [0]
01 ADD 1
02 STA [0]
03 HLT
\end{slsourcecode}
\ankinote{syntax}{Syntax}{Befehle, Zahl, [Adresse]}
Es gibt hier keine höheren Konzepte wie eingebaute Schleifen; Schleifen müssen selbst entwickelt werden, indem ein konditionaler Jump-Befehl genutzt wird um Befehle zu wiederholen.

\subsection{Arbeitsregister}
Der Akkumulator speichert zwischenzeitig Daten, wie Operanten oder Ergebnisse. Der Befehlszähler speichert die Adresse zur auszuführenden Befehls.
\ankinote{arbeitsregister}{Arbeitsregister}{Akkumulator (Werte speichern), Befehlszähler}

\subsection{Ausführung}
Zuerst wird der bevorstehende Befehl aus dem Arbeitsspeicher geladen und dekodiert, bei Bedarf werden noch Operanten geladen. Der Befehl wird ausgeführt. Das Befehlszählregister wird um eine erhöht und der Prozess fängt von vorne an, mit dem nächsten Befehl.

\end{document}